\chapter{Architektura i struktura aplikacji}

\section{Architektura aplikacji}

Aplikacja została przygotowana jako program desktopowy WPF. Projekt wykorzystuje prostą architekturę warstwową z elementami wzorca MVVM. Interfejs użytkownika znajduje się w pliku XAML, logika prezentacji w klasie \texttt{MainViewModel}, a operacje biznesowe w serwisie \texttt{HotelService}.

\begin{figure}[H]
\centering
\includegraphics[width=0.95\textwidth]{figures/Architektura.eps}
\caption{Uproszczona architektura aplikacji HotelManagementApp}
\end{figure}

\section{Podział projektu na warstwy i moduły}

Projekt można podzielić na kilka części:

\begin{itemize}
  \item warstwa modeli, czyli klasy opisujące dane w systemie,
  \item warstwa danych, czyli klasa \texttt{HotelDbContext} oraz migracje,
  \item warstwa repozytoriów, czyli klasy \texttt{IRepository} i \texttt{Repository},
  \item warstwa usług, czyli \texttt{IHotelService} i \texttt{HotelService},
  \item warstwa ViewModel, czyli \texttt{MainViewModel},
  \item warstwa interfejsu, czyli \texttt{MainWindow.xaml} oraz \texttt{MainWindow.xaml.cs}.
\end{itemize}

\section{Struktura folderów oraz organizacja kodu}

Struktura projektu została podzielona według odpowiedzialności plików.

\begin{lstlisting}[style=textStyle,caption={Uproszczona struktura folderów projektu},label={lst:struktura}]
HotelManagementApp
|-- Commands
|   |-- RelayCommand.cs
|-- Data
|   |-- HotelDbContext.cs
|   |-- DbSeeder.cs
|-- Migrations
|   |-- 20260608000000_InitialCreate.cs
|   |--HotelDbContextModelSnapshot.cs
|-- Models
|   |-- Guest.cs
|   |-- Room.cs
|   |-- RoomType.cs
|   |-- Reservation.cs
|   |-- Payment.cs
|   |-- Employee.cs
|   |-- Service.cs
|   |-- ReservationService.cs
|-- Repositories
|   |-- IRepository.cs
|   |-- Repository.cs
|-- Services
|   |-- IHotelService.cs
|   |-- HotelService.cs
|-- ViewModels
|   |-- MainViewModel.cs
|   |-- ObservableObject.cs
|-- App.xaml.cs
|-- App.xaml
|-- MainWindow.xaml
|-- MainWindow.xaml.cs
|-- HotelManagementApp.csproj
\end{lstlisting}

Folder \texttt{Models} zawiera klasy mapowane na tabele w bazie danych. Folder \texttt{Data} zawiera konfigurację bazy danych i dane startowe. Folder \texttt{Services} zawiera logikę aplikacji, a folder \texttt{ViewModels} zawiera dane i komendy używane przez interfejs WPF.

\section{Zastosowane wzorce projektowe}

W projekcie zastosowano kilka prostych wzorców i mechanizmów organizacji kodu. Najważniejszym z nich jest MVVM, ponieważ widok korzysta z danych i komend udostępnionych przez \texttt{MainViewModel}. Zastosowano również Dependency Injection w pliku \texttt{App.xaml.cs}, dzięki czemu klasy nie tworzą ręcznie swoich zależności.

Dostęp do danych został częściowo oddzielony za pomocą wzorca Repository. Logika biznesowa znajduje się w serwisie \texttt{HotelService}. Dzięki temu kod odpowiedzialny za interfejs nie musi bezpośrednio wykonywać wszystkich operacji na bazie danych.

\section{Mechanizmy walidacji danych}

Walidacja danych została wykonana w serwisie. Program sprawdza między innymi, czy imię, nazwisko, numer telefonu i adres e-mail gościa nie są puste oraz czy adres e-mail zawiera znak \texttt{@}. Przy rezerwacji sprawdzane jest, czy data wymeldowania jest późniejsza niż data zameldowania oraz czy pokój jest dostępny w podanym terminie.

Dla usług hotelowych sprawdzana jest nazwa i cena. Nazwa nie może być pusta ani zbyt długa, a cena nie może być ujemna. Dzięki temu użytkownik nie może zapisać podstawowo błędnych danych.

\section{Obsługa błędów i wyjątków}

W projekcie zastosowano prostą obsługę błędów. Metody ViewModelu wywołują operacje serwisu wewnątrz bloków \texttt{try/catch}. Jeżeli wystąpi błąd, komunikat zostaje pokazany użytkownikowi w dolnej części okna.

Przykładowo, gdy użytkownik wybierze datę wymeldowania wcześniejszą niż data zameldowania, system nie zapisze rezerwacji i pokaże informację o błędzie. Podobnie aplikacja reaguje na próbę dodania płatności z kwotą mniejszą lub równą zero.

\section{Czytelność i organizacja projektu}

Nazwy klas odpowiadają elementom domeny hotelowej, na przykład \texttt{Guest}, \texttt{Room}, \texttt{Reservation}, \texttt{Payment} i \texttt{Service}. Dzięki temu łatwo zrozumieć, za co odpowiada dana część kodu. Projekt nie jest nadmiernie skomplikowany, ale spełnia wymagania dotyczące podziału na warstwy, obsługi bazy danych i organizacji interfejsu użytkownika.
